\documentclass[10pt,a4paper]{amsart}
\usepackage[margin=19mm]{geometry}
\usepackage{amsmath,amssymb,mathtools}
\usepackage[scr=rsfs,cal=euler]{mathalfa}
\usepackage{xfrac}
\usepackage[unicode,hidelinks,bookmarks=false]{hyperref}
\newcommand{\ie}{i.e.\ }
\newcommand{\cf}{cf.\ }
\newcommand{\resp}{respectively\ }
\newcommand{\Subsection}{\subsection}
\providecommand{\nobreakdash}{\nobreak\hbox{-}}
\setlength{\emergencystretch}{2em}
\multlinegap0pt
\usepackage{bm}
\usepackage{bbm}
\usepackage{dsfont}
\usepackage{xspace}
\usepackage{cancel}
\usepackage{mathtools}
\usepackage{amsthm}
\makeatletter
\@ifundefined{theorem}{\newtheorem{theorem}{Theorem}}{}
\@ifundefined{lemma}{\newtheorem{lemma}[theorem]{Lemma}}{}
\makeatother
\newcommand{\CircuitVal}{\text{CircuitVal}}
\newcommand{\PV}{\mathsf{PV}}
\newcommand{\UP}{\mathsf{UP}}
\newcommand{\abs}[1]{\lvert #1 \rvert}
\newcommand{\bd}{\mathsf{b}}
\newcommand{\poly}{\mathsf{poly}}
\title[Parallelism and Adaptivity in Student-Teacher Witnessing]{Parallelism and Adaptivity in Student-Teacher Witnessing}
\author{Ond\v{r}ej Je\v{z}il, Dimitrios Tsintsilidas}
\date{Source: arXiv:2602.19934v1 (CC BY 4.0)}
\begin{document}
\maketitle

\noindent\emph{Source and licence.} Parallelism and Adaptivity in Student-Teacher Witnessing. Version arXiv:2602.19934v1 (CC BY 4.0), distributed under the Creative Commons Attribution 4.0 International licence, as recorded in the arXiv licence metadata for this exact version and shown on the abstract page https://arxiv.org/abs/2602.19934v1. Full rights notice: licenses/arxiv-2602.19934-CC-BY-4.0.txt.\par\smallskip

\noindent\textit{Editorial scope.} Counted unit: the Theorem whose source label is \texttt{theoremkobb} in the article (source0.tex lines 1181--1184, source label \texttt{theoremkobb}), delivered together with the article's own complete original proof of it (source0.tex lines 1185--1212, one \texttt{proof} environment of 28 lines). No mathematics is written, repaired or re-derived: statement and proof are the article's own text line for line. Required context retained uncounted: thm:bbwitness (lines 957--959); lemmarefuter (lines 1172--1175). Every definition, display and label of those lines is retained unchanged. Omitted from this package and not redistributed: 1--956, 960--1171, 1176--1180, 1213--1383. Nothing inside a retained excerpt is omitted or altered: every retained body is delivered exactly as the article's lines are written, so no omission receipt of any kind accompanies this package. Citations: the retained excerpts carry no citation command at all, so no bibliography entry of the article is transcribed and no reference is redistributed. Reference adaptation: none was needed and none was made; every reference inside the delivered unit resolves inside it, so no unresolved reference or citation is produced. Printed slips kept literal: no printed word, symbol, hypothesis or number of the counted statement or of its proof was altered. Layout and class: the article's own class is \texttt{article} with packages including graphicx, fullpage, amsfonts, indentfirst, amsmath, amsthm, amssymb, natbib, hyperref, cleveref, color, algorithm2e, enumitem, tikz; the wrapper is the lane's pinned \texttt{amsart} editorial wrapper at 10pt on A4, which restates the theorem-like environments the retained text needs and adds the article's own macro definitions that the retained text needs (\texttt{\textbackslash CircuitVal}, \texttt{\textbackslash PV}, \texttt{\textbackslash UP}, \texttt{\textbackslash abs}, \texttt{\textbackslash bd}, \texttt{\textbackslash poly}), with extra wrapper packages (bm, bbm, dsfont, xspace, cancel, mathtools). The delivered numbering is the unit's own. Assets: no retained span contains a figure environment, a graphics-inclusion command, a PDF-page inclusion, a table or a TikZ picture; no image file is copied into this package, the package declares no asset dependency and no image file is redistributed with it. Licence: version arXiv:2602.19934v1 of this article is distributed under the Creative Commons Attribution 4.0 International licence, as recorded in the arXiv licence metadata for this exact version and shown on the abstract page https://arxiv.org/abs/2602.19934v1. A full rights notice is delivered at licenses/arxiv-2602.19934-CC-BY-4.0.txt. Characters: every retained excerpt is pure ASCII, so the delivered engine, XeLaTeX with its default Latin Modern text font, reports no missing character.
\begin{theorem} \label{thm:bbwitness}
    Let $\PV_j+\mathsf{BB}(\Sigma^b_{j},\bd)\vdash (\forall x)(\exists y)(\forall z\leq t)(\varphi(x,y,z))$, where $\varphi$ is open and $t$ is a $\PV_j$-term, then the $\mathbf{TF\Sigma^p_{j+1}}$ problem corresponding to $\varphi$ is in $\mathcal{ST}^{\Sigma^p_{j-1}}[O(1),\poly(\bd)]$.
\end{theorem}

\begin{lemma}\label{lemmarefuter}
    Let $f,\tilde f, g$ and $p$ be as in the previous lemma. Then there exists $c_f$ and a $\PV$-symbol $e$ such that 
    \[\PV_1+\mathsf{BB}(\Sigma^b_1)\vdash (\forall 1^{(n)},n\geq c_f)(\forall i\leq \abs{p(1^{(n)})})(g(e(i,1^{(n)}))\neq \CircuitVal(f_i(1^{(n)}),e(i,1^{(n)}))).\]
\end{lemma}

\begin{theorem}\label{theoremkobb}
    For every $k\geq 1$ there is a unary $\PV$-symbol $h$ such that for every constant $c\geq 1$
    \[\PV_1+\mathsf{BB}(\Sigma^b_1)\nvdash \UP_{k,c}(h).\]
\end{theorem}

\begin{proof}
    We will proceed by contradiction. Assume that there is a $k$ where for any unary $h$ there is a constant $c$ such that \[\PV_1+\mathsf{BB}(\Sigma^b_1)\vdash \UP_{k,c}(h).\]

    In particular, there is a constant $c\geq1$ such that $\PV_1+\mathsf{BB}(\Sigma^b_1)\vdash \UP_{k,c}(g).$ By Theorem~\ref{thm:bbwitness}, there is a polynomial-time student $s$ which generates sets of polynomially many answers for the candidate circuit computing the symbol $g$ and after $r$-many rounds of correction from the teacher it always outputs at least one circuit which computes $g$ on the input length $1^{(n)}$. 
    
    Assume that $s^1,\dots,s^{r}$ are the $\PV$-symbols computing the set of answers of the student from the teachers answer for each of the $r$-many rounds, these can be straightforwardly constructed from the subterms of $s$. That is, there is a $\PV$-symbol $p$ such that
    \begin{align*}
        \mathbb{N}\models \quad&(\exists i\leq \abs{p(1^{(n)})})(g((x_1)_i)=\CircuitVal(s^1_i(1^{(n)}),(x_1)_i))\\
                          \lor&(\exists i\leq \abs{p(1^{(n)})})(g((x_2)_i)=\CircuitVal(s^2_i(1^{(n)},x_1),(x_2)_i))\\
                          &\quad\vdots\\
                          \lor&(\exists i\leq \abs{p(1^{(n)})})(g((x_r)_i)=\CircuitVal(s^r_i(1^{(n)},x_1,\dots,x_{r-1}),(x_{r})_i)),
    \end{align*}
    where $(-)_i$ denotes the $\PV$-symbol which extracts the $i$-th element of a sequence. In the rest of the proof we will show that from the assumption, that there is a $\tilde c_1$ such that \[\PV_1+\mathsf{BB}(\Sigma^b_1)\vdash \UP_{k,\tilde c_1}(s^1(1^{(n)})),\] we can invoke Lemma~\ref{lemmarefuter} to falsify the first disjunct. Repeating this $r$-many times, we obtain a contradiction.
    
    By our assumption, there is a $\tilde c_1\geq 1$ such that $\PV_1+\mathsf{BB}(\Sigma^b_1)\vdash \UP_{k,\tilde c_1}(s^1(1^{(n)}))$, by Lemma~\ref{lemmarefuter} there is a $\PV$-symbol $e^1$ and $c_{1}$, such that
    \[\PV_1+\mathsf{BB}(\Sigma^b_1)\vdash (\forall 1^{(n)},n\geq c_1)(\forall i\leq \abs{p(1^{(n)})})(g(e^1(i,1^{(n)}))\neq\CircuitVal(s^1_i(1^{(n)}),e^1(i,1^{(n)}))),\]
    by substituting $\tilde e^1(1^{(n)})$ for $x_1$, where $\tilde e^1(1^{(n)})$ is the $\PV$-symbol computing the sequence \[(e^1(i,1^{(n)}))_{i\leq \abs{p(1^{(n)})}},\] we falsify the first disjunct and obtain
    \begin{align*}
        \mathbb{N}\models \quad&(\exists i\leq \abs{p(1^{(n)})})(g((x_2)_i)=\CircuitVal(s^2_i(1^{(n)},\tilde e^1(1^{(n)})),(x_2)_i))\\
                          \lor&(\exists i\leq \abs{p(1^{(n)})})(g((x_3)_i)=\CircuitVal(s^3_i(1^{(n)},\tilde e^1(1^{(n)}),x_2),(x_3)_i))\\
                          &\quad\vdots\\
                          \lor&(\exists i\leq \abs{p(1^{(n)})})(g((x_r)_i)=\CircuitVal(s^r_i(1^{(n)},\tilde e^1(1^{(n)}),x_2,\dots,x_{r-1}),(x_{r})_i)).
    \end{align*}

    To proceed further, we again use our assumption to obtain $\tilde c_2\geq 1$ such that 
    \[\PV_1+\mathsf{BB}(\Sigma^b_1)\vdash \UP_{k,\tilde c_2}(s^2(1^{(n)},\tilde e^1(1^{(n)}))),\]
    and apply Lemma~\ref{lemmarefuter} again. After $r$-many invokation of Lemma~\ref{lemmarefuter}, we obtain that the empty disjunction is true in $\mathbb{N}$ which is a contradiction.
\end{proof}

\end{document}
